
\section{\method}
\label{sec:framework}


This section presents the design of \method,
a unified framework for long-term life simulation in agent societies.
The framework is organized around three core concepts: the simulation procedure, the agent, and the environment model:
\textit{(1)} For the simulation procedure, agents live through weekly cycles, each consisting of four stages: Plan, Contact, Activity, and Review, which jointly support diverse social behaviors.
\textit{(2)} For the agent, the core challenge is context management, \ie, providing each LLM with sufficient context including the character's persona, states, relationships, and memories. 
Agents' long-term memory is implemented as a file system that agents autonomously manage via function calls. 
\textit{(3)} To orchestrate the simulation, \method introduces the environment model, a separate LLM that organizes events, provides environmental feedback, and drives the simulation forward, replacing hard-coded rules.
We describe agent design in \S~\ref{sec:agent_design},
the simulation procedure in \S~\ref{sec:simulation_procedure},
contact and scheduling in \S~\ref{sec:contact},
activities in \S~\ref{sec:activity},
and environment design in \S~\ref{sec:environment}.

\subsection{Agent Design}
\label{sec:agent_design}

In \method, each agent role-plays as a specific persona, which comprises a \textit{profile}, \textit{social relationships}, and \textit{dynamic states}, comprehensively representing the character.


\begin{enumerate}
    \item \textbf{Profile}:
    The profile stores an agent's identity information, including a background story, personality traits, talents, as well as initial values for positions and assets. 
    These attributes remain relatively stable during simulation and are only updated once per simulated year.
    The full profile schema is provided in \S~\ref{app:schema}, and the character creation process is described in \S~\ref{app:creation}.
    
    \item \textbf{Social Relationships}:
    In \method, relationships are modeled entirely via inter-character memory.
    We do not explicitly define or store character relationships.
    Instead, relationships are represented through each character's memory of the other, stored as free text and autonomously managed by agents.
    This design uniformly represents friends, lovers, strangers, rivals, and other relationship types without requiring separate mechanisms.
    Relationships can be unidirectional and expanded through multi-person activities.
    The relationship network is initialized during world creation (\S~\ref{app:creation}) and subsequently managed autonomously by agents.
    
    \item \textbf{Dynamic States}:
    Dynamic states represent an agent's current life conditions and evolve with activities.
    They include vitality, fulfillment, skills, position, and assets.
    \textit{Vitality} represents energy level. 
    \textit{Fulfillment} is based on Maslow's hierarchy of needs~\citep{maslow1943theory}, which captures agents' subjective satisfaction from four dimensions: mood, material, social, and esteem.
    It serves as the basis for subjective rewards.
    Following hedonic adaptation theory~\citep{Headey2014}, fulfillment decays naturally each week.
    Fulfillment and vitality changes are determined by the environment model based on activity content, rather than hard-coded rules.
    The details are provided in \S~\ref{sec:appendix_states}.
    \textit{Skills} represent an agent's acquired abilities.
    \textit{Position} represents an agent's current occupation, which provides weekly income and skill growth (\S~\ref{sec:appendix_enrollment}).
    \textit{Assets} consist of a deposit and a list of possessions (\S~\ref{app:economy}).
    
\end{enumerate}

Besides, \method carefully designs a context management mechanism to provide agents with sufficient context across stages, and equips agents with memory files as long-term memory that agents autonomously manage.
Together, \method provides rich context for LLMs to simulate human social lives.

\begin{enumerate}
    
    \item \textbf{Context Management}:
    \method carefully designs a comprehensive context management mechanism to provide agents' underlying LLMs with sufficient information to simulate human social lives.
    This is achieved through three context layers:
    \textit{(1)} \textbf{Roleplay prompt}: provides the foundational role-playing information shared across all stages.
    It contains the agent's full persona, recent weekly diaries as short-term memory, summaries of key memory files as long-term memory, worldview rules and role-play principles, \etc,
    enabling agents to correctly accumulate memory across stages and weeks throughout the simulation.
    \textit{(2)} \textbf{Stage prompt}: provides stage-specific instructions and context, \eg, scheduling rules and communication history  for the Contact stage.
    \textit{(3)} \textbf{Message history}: records LLMs' multi-turn messages accumulated within the current stage, including prior dialogues, calls and results of memory functions, and compacted reasoning process.
    Together, these three layers combine automatic and agent-driven context management, providing comprehensive context for agents to simulate human social lives.
    The details are elaborated in \S~\ref{sec:appendix_context}.

    \item \textbf{Memory Files}:
    While weekly diaries serve as short-term memory by recording recent experiences, memory files provide agents with \textit{long-term memory}, a file-system-based persistent store that agents autonomously manage via function calls.
    Agents decide what to remember, update, or discard, giving them full control over what persists across weeks and years.
    Each agent maintains memory files in three categories:
    \texttt{general.txt} for personal notes and plans,
    \texttt{characters/<who>.txt} for the agent's knowledge of and relationship with specific people,
    and \texttt{others/<name>.txt} for any other topics.
    Memory files provide context to LLMs in two ways: (1)
    summaries of recently accessed files are \textit{automatically} included in the roleplay prompt, 
    and (2) agents \textit{autonomously} retrieve full content on demand, included in the message history.
    Agents manage these files through three functions: \texttt{read\_file} to retrieve content, \texttt{update\_file} to write new content, and \texttt{list\_files} to list all files.
    A read-before-write constraint is enforced: an update is only permitted after the agent has read the target file in the same invocation, ensuring updates build upon existing content rather than overwriting blindly.
    Agents can also create new files for topics they wish to track.
    
    
    \end{enumerate}
\subsection{Simulation Procedure}
\label{sec:simulation_procedure}

A fundamental constraint of LLM-based agents is that they generate responses turn by turn:
they receive a complete context and produce a single reply, unlike humans who continuously perceive and react in real time (\eg, interrupting a conversation mid-sentence).
The simulation procedure must therefore be designed around turn-based interaction.
Meanwhile, \method focuses on abstract social interactions (\eg, planning, socializing, decision-making) rather than low-level operations (\eg, movement, object manipulation).
This abstraction enables the system to model richer social behaviors with fewer LLM calls, achieving higher token efficiency for long-term simulation.

With these considerations, we design a comprehensive simulation procedure.
\method uses the \textit{week} as its basic time unit.
A simulated year spans $n_w$ weeks; year-end triggers profile updates, position application, and life reward calculation.~\footnote{%
The \textit{day}, \textit{week}, and \textit{year} in \method serve as three hierarchical layers of abstract simulation time.
They are designed to simulate human life cycles at comparable scales, without strictly corresponding to real-world calendar units.}
Each week consists of four stages: Plan, Contact, Activity, and Review.




\begin{enumerate}

\item \textbf{Plan}: 
Based on memory and current state, each agent makes a weekly plan.
Agents are able to review previous plans and store or update new ones through function calls (\eg, \texttt{read\_file}, \texttt{update\_file}).
Agents also select an abstract consumption level (\ie, living standard) for the week, which abstracts their monetary spending and influences material fulfillment.

\item \textbf{Contact (\& Scheduling)}: 
In this stage, agents communicate with each other over multiple rounds.
One primary purpose of communication is to arrange shared schedules: agents can propose invitations, which others may accept or decline, thereby creating joint activity schedules.
The detailed mechanism is introduced in \S~\ref{sec:contact}.

\item \textbf{Activity}: 
In this stage, agents carry out diverse activities.
Each weekly cycle contains $n_d$ active days, and on each day, every agent carries out exactly one activity.
Activities fall into four types: joint, solo, encounter, and public.
Through these activities, agents can pursue personal growth, work, or leisure on their own, or engage with others in shared experiences.
Activity types and execution are detailed in \S~\ref{sec:activity}.

\item \textbf{Review}: 
At the end of each week, agents reflect on their weekly experiences and summarize them into a weekly diary.
Based on this reflection, agents can update their memory files via function calls.
The weekly diaries are included in each agent's context in subsequent weeks, allowing agents to recall their past experiences.

\end{enumerate}

At the end of each simulated year, three processes take place.
\textit{(1) Profile Update:} the environment model makes incremental updates to each agent's profile (\eg, personality traits, talents) based on their accumulated experiences over the year.
\textit{(2) Position Application:} agents may apply for new positions based on their preferences and compete for limited vacancy based on their abilities and skills.
\textit{(3) Life Reward Calculation:} life rewards are calculated, measuring each agent's social standing, subjective fulfillment, and economic status over the past year.

\subsection{Contact \& Scheduling}
\label{sec:contact}

The Contact stage serves two purposes: communication between agents and scheduling of activities.
Communication is restricted to pairwise exchanges, while multi-person conversations are deferred to joint activities.
Constrained by the turn-based nature of LLM generation, communication in \method proceeds in rounds.
Each week consists of $n_c$ contact rounds.
In each round, an agent receives newly arrived messages along with the contact history from the past three weeks, then decides whom to contact and what actions to take.
After all rounds conclude, the system resolves the contact history to determine which joint activities have been successfully created.

\paragraph{Actions}
In \method, agents must generate \textit{actions} to trigger real communication and scheduling.
Actions are wrapped with \texttt{<role\_action>} tags in agents' responses, which the system then parses and executes.
They are different from function calls.
Four action types are available: \texttt{contact}, \texttt{propose\_joint\_activity}, \texttt{respond\_invitation} and \texttt{cancel\_joint\_activity} (\S~\ref{sec:appendix_contact_details}).

\paragraph{Schedule Resolution} After all contact rounds conclude, the system resolves which joint activities are successfully created based on all collected actions.
The process handles cancellations, deduplicates responses, resolves time conflicts, and checks creation condition. 
The full rules and procedure are described in \S~\ref{sec:appendix_procedure}.

\paragraph{Scheduling Public and Encounter Activities}
Public and encounter activities are scheduled in dedicated stages before and after Contact, respectively.
Before Contact, the environment model creates public events for the upcoming weeks, and agents sign up for those that match their interests.
After Contact, the environment model arranges encounter activities to create chance meetings for idle agents, \ie, those without any scheduled activity on a given day.

\subsection{Activities}
\label{sec:activity}

Activities are the core stage of \method, where agents carry out activities and develop abilities, fulfillment, wealth, and social relationships.
\method supports four activity types: joint, solo, encounter, and public.
Agents carry out activities according to their schedules, defaulting to solo activities when unscheduled.
After each activity, the environment model provides feedback based on the agent's experience (including state changes), and the agent reflects on the outcome and may update its memory files accordingly.
The four activity types are described below:

\begin{enumerate}
\item \textbf{Joint Activity}:
Joint activities model multi-agent, multi-turn interactions.
They are created through invitation and negotiation during the Contact stage.
Participants take turns speaking; following~\citet{wang2025coser}, we employ the environment model to provide environmental descriptions and select the next speaker each turn.
Agents can choose the visibility of their generated content, using special tags to mark each segment as public, private, or selective (visible only to specified persons).
Agents are also permitted two actions: \texttt{gift} to transfer items to another participant, and \texttt{exit} to leave the activity early.
Additionally, we introduce a response filtering mechanism, where the environment model evaluates each agent response based on a set of roleplay principles covering anthropomorphism, character fidelity, and feasibility, filtering out responses that violate any principle.
Details about joint activity execution and response filtering are provided in \S~\ref{sec:appendix_activity} and \S~\ref{sec:appendix_principles}, respectively.

\item \textbf{Solo Activity}:
Solo activity serves as the default when an agent has no other schedule, allowing activities such as studying, working, or leisure and consumption.
It follows a single-turn format: the agent describes its intended action, and the environment model evaluates the feasibility of the action based on the character's background and provides an outcome.
If the agent's intention does not match its abilities, the environment model provides negative feedback.
Solo activities also support spending, where agents may purchase goods or services to gain material fulfillment (\S~\ref{app:economy}).
Details about solo activity execution are described in \S~\ref{sec:appendix_solo_activity}.

\item \textbf{Encounter Activity}:
Encounters are chance meetings arranged by the environment model for idle agents, and serve as a special case of joint activity.
They follow the same multi-turn dialogue format as joint activities, but do not appear in agent schedules.
Encounters are designed to let the environment model create meaningful chance meetings: for example, introducing strangers to expand social networks, or bringing together agents with significant plot connections or special relationships.

\item \textbf{Public Activity}:
Public activities represent open community events centered around shared interests (\eg, workshops, interest groups).
They are designed by the environment model, and agents may sign up independently based on their own interests.
The execution generally follows a similar flow to solo activity.
The key difference is that, at the end of the activity, agents see what other participants did, and may get to know them.
Details for public activities are described in \S~\ref{sec:appendix_public_activity}.

\end{enumerate}

\subsection{Environment Design}
\label{sec:environment}

This section introduces the environment design in \method, including the environment model, the economy system, the position application mechanism, and the location system.
These components serve as essential parts that support the comprehensive simulation of social lives in agent societies.

\paragraph{Environment Model}
Building a comprehensive social simulation environment to appropriately orchestrate the simulation and provide environmental feedback for role-playing agents is a complex yet important task.
In \method, we introduce the environment model, a stateless LLM, to handle these functions instead of hard-coding extensive rules.
During activities, it provides feedback to agent behavior, judging feasibility and evaluating outcomes.
Besides, it predicts next speakers in joint activities, generates public and encounter activities, ranks candidates for position applications, updates agent profiles at year-end, and performs response filtering to exclude low-quality outputs.
These functions collectively support the smooth progression of life simulation in the multi-agent society.
Hence, the environment model essentially serves as a generative environment engine, providing intelligent, generative responses to agent behaviors.
Our prompts for the environment model are displayed in Tables~\ref{tab:prompts_god_joint}, \ref{tab:prompts_god_eval}, and~\ref{tab:prompts_god_world}.

\paragraph{Economy System}
\method includes a basic economy system where agents earn and spend money.
Income comes from three sources: a position-based weekly salary, character-specific weekly income, and additional work during the activity stage.
Expenses take two forms.
First, agents select a living standard each week (from frugal to luxurious), representing the richness of their material life for that week.
Second, agents may choose to engage in consumption activities during solo activities (\eg, shopping, entertainment).
Through spending, agents gain material fulfillment, which contributes to improvement in subjective reward.
Meanwhile, accumulating more wealth leads to higher economy reward (\S~\ref{app:economy}).

\paragraph{Position System}
Each agent holds a position that represents its job or social role.
A position provides two weekly benefits: a fixed income and skill growth in relevant abilities.
During character construction, the environment model assigns each agent an initial position based on the world setting and character profiles.
Once per year, a position application process allows agents to decide whether to change their position, and the environment model determines outcomes based on agents' abilities and position requirements.
The details about position and position application are described in \S~\ref{sec:appendix_enrollment}.

\paragraph{Location System}
\method implements a simple location mechanism.
Its primary purpose is to provide agents with grounded environmental perception.
In our early experiments, we find that agents tend to hallucinate environmental details and objects without location information explicitly provided.
Hence, we employ the environment model to create a set of locations for each world, and require joint activities and encounter activities to specify a location.
Details are provided in \S~\ref{sec:appendix_location}.
